
\lesson{03/11}
\textbf{Proof of Wick's theorem (by induction)}\\
The base case is with two fields and it is easily proven by brute force.

Supposing the theorem holds for $n$ fields, consider the time-ordered product
of $n+1$ fields: without losing generality, consider $x^0_{n+1}<x_i^0$, then
$$T\left\{\phi(x_1)...\phi(x_{n+1})\right\}=T\left\{\phi(x_1)...\phi(x_n)
\right\}\phi(x_{n+1})$$
from the hypothesis, Wick's theorem can be used
$$T\left\{\phi(x_1)...\phi(x_{n+1})\right\} =\left( N\left\{\phi(x_1)...\phi(x_n)
\right\} + \sum_{}^{}\text{all contractions}\right)\phi(x_{n+1}).$$
Consider now
$$N\left\{\phi(x_1)...\phi(x_n)\right\}\left(\phi(x_{n+1})^++\phi(x_{n+1})^-
\right).$$
Noting that the normal-ordered product can be written as
$$N\left\{\phi(x_1)...\phi(x_n)\right\} = \sum_{A,B}\prod_{i\in A}\phi^+(x_i)
\prod_{j\in B}\phi^-(x_j)$$
(where $A,B$ are partitions of the set of indexes) yields
\begin{gather*}
N\left\{\phi(x_1)...\phi(x_n)\right\}\phi(x_{n+1}) = 
\sum_{A,B}\prod_{i\in A}\phi^+(x_i)\prod_{j\in B}\phi^-(x_j)\left(\phi^+(x_{n+1})
+\phi^-(x_{n+1})\right).
\end{gather*}
The $\phi^-(x_{n+1})$ is not problematic, since it is already normal ordered, but
in order to express this normal-ordered product in terms of $N\left\{\phi(x_1)...
\phi(x_{n+1})\right\}$, the $\phi^+(x_{n+1})$ needs to go through all the
$\phi^-(x_j)$'s yielding commutators:
\begin{gather*}
N\left\{\phi(x_1)...\phi(x_n)\right\}\phi(x_{n+1}) = 
N\left\{\phi(x_1)...\phi(x_{n+1})\right\} + \sum_{A,B}^{}\prod_{i\in A}\phi^+(x_i)
\phi^+(x_{n+1})\prod_{j\in B}\phi^-(x_j) +\\+\sum_{A,B}\prod_{k\in B}\phi^+(x_i)
\prod_{j\neq k \in B} \phi^-(x_j)\overline{\phi(x_k)\phi(x_{n+1})}
\end{gather*}
So the contractions sum with the others and the normal-ordered product now
includes $\phi(x_{n+1})$ \qed

\textbf{Example: four fields}\\
Denoting $\phi_i = \phi(x_i)$, consider
$$T\left\{\phi_1\phi_2\phi_3\phi_4\right\} = N\left\{
    \overline{\phi_1\phi_2}\phi_3\phi_4 + \overline{\phi_2\phi_3}\phi_1\phi_4 +
    \overline{\phi_3\phi_4}\phi_1\phi_2 + \overline{\phi_1\phi_4}\phi_2\phi_3+
    \overline{\phi_1\phi_2}\;\overline{\phi_3\phi_4}+
    \overline{\phi_1\phi_4}\;\overline{\phi_2\phi_3}+
    \overline{\phi_1\phi_3}\;\overline{\phi_2\phi_4}
\right\}$$
When taking the vacuum expectation value the first 4 terms vanish because of $N$:
$$\bra{0}T\left\{\phi_1\phi_2\phi_3\phi_4\right\}\ket{0}=
i\Delta_F(x_1-x_2)i\Delta_F(x_3-x_4)+
i\Delta_F(x_1-x_4)i\Delta_F(x_2-x_3)+
i\Delta_F(x_1-x_3)i\Delta_F(x_2-x_4)
$$
\textbf{Note} The proof for Dirac fields is similar, save a $(-1)^p$ sign, given
by the fact that $p$ anticommutation relations need to be used.
\subsection{The S-matrix}
This is the greatest success of Quantum Field Theory: the predictive part
of the Standard Model has passed with flying colours every test yet devised;
the core of the predictions is the $S$-matrix.

As said before, the $S$-matrix is defined as
$$\lim\limits_{T\to\infty}U_I(T,-T)=\lim\limits_{T\to\infty}
T\left\{\exp\left[-i\int_{-T}^{T} \dd t H_I(t) \right]\right\}
:=S$$
Since the exponential is just notation for its series, the $S$-matrix can also
be written as
$$S = T\left\{\mathbb{I} -i\int \dd^4x \mathcal{H}_I(x^\mu) + \frac{(-i)^2}
{2!}\int\dd^4y\int_{}^{} \dd^4w \mathcal{H}_I(y^\mu)\mathcal{H}_I(w^\mu) 
\right\}$$
where $\mathcal{H}_I$ is the interacting picture Hamiltonian density.\\
The reason behind the $T\to\pm\infty$ is to study \textbf{scattering processes}
by considering the particles as non-interacting in the initial and final state:
e.g. the initial state may be
$$a^\dagger(\mathbf{k})a^\dagger(\mathbf{k'})\ket{0}$$
which is a state with $2$ scalar \textit{free} particles (as $T\to-\infty$) of
momenta $\mathbf{k,k'}$, while the final state may be
$$\bra{0}a(\mathbf{k'})a(\mathbf{k})$$
which are again \textit{free} particles (as $T\to\infty$): the probability
amplitude of this happening is
$$A=\bra{\text{final}}S\ket{\text{initial}}:$$
the convenient thing is that the $S$-matrix can be computed perturbatively using
the exponential series.\\
\textbf{Note} The order of the series is the \textit{number of vertices} in a
\textbf{Feynman Diagram}.
\subsection{The $\phi^4$ self-interaction}
The simplest interaction Hamiltonian density is
$$\mathcal{H}_I = \frac{\Lambda}{4!}\phi^4(x)$$
where $\Lambda$ is a parameter.

